\documentclass[article]{ajs}

\usepackage{amsmath}

%%%%%%%%%%%%%%%%%%%%%%%%%%%%%%
%% declarations for ajs.cls %%
%%%%%%%%%%%%%%%%%%%%%%%%%%%%%%

\author{Arquímedes De León Chacón Chacón\,\orcidlink{0000-0002-7014-7513}\\ iMetric, Spain}

%% WORKING TITLE -- candidates discussed, easy to swap later:
%%  1. Stress-Testing Chatterjee's Xi Correlation Coefficient: A Monte Carlo
%%     Study of Sample Size, Distribution Shape, and Variable Type
%%  2. A Factorial Monte Carlo Evaluation of Chatterjee's Xi Correlation
%%     Coefficient Against Established Measures of Dependence
%%  3. Chatterjee's Xi Correlation Coefficient Under Stress: Sample Size,
%%     Distribution Shape, and Variable Type Effects in a Monte Carlo
%%     Framework
\title{Stress-Testing Chatterjee's Xi Correlation Coefficient: A Monte Carlo Study of Sample Size, Distribution Shape, and Variable Type}

\Plainauthor{Arquímedes De León Chacón Chacón}
\Plaintitle{Stress-Testing Chatterjee's Xi Correlation Coefficient: A Monte Carlo Study of Sample Size, Distribution Shape, and Variable Type}
\Shorttitle{Stress-Testing Chatterjee's Xi Correlation Coefficient}

\Abstract{
  Chatterjee's $\xi$ correlation coefficient (Chatterjee, 2021) was
  proposed as a simple, rank-based measure capable of detecting arbitrary
  functional dependence between two variables, together with a
  computationally cheap asymptotic test of independence. Since its
  introduction, a fast-growing theoretical literature has proposed
  refinements, bias corrections, and multivariate extensions (e.g.,
  Azadkia and Chatterjee, 2021), each typically accompanied by a small,
  narrowly scoped simulation. No published study has stress-tested
  $\xi$ across a full factorial design crossing sample size,
  distribution shape, and variable type, nor benchmarked it
  systematically against established nonlinear dependence measures such
  as distance correlation and copula-based statistics under the same
  conditions. \emph{[Placeholder --- to be completed once the
  simulation design and results are finalized.]}
}

\Keywords{chatterjee's xi, rank correlation, monte carlo simulation, measures of dependence, distance correlation, robustness, \proglang{R}}
\Plainkeywords{chatterjee's xi, rank correlation, monte carlo simulation, measures of dependence, distance correlation, robustness, R}

%% publication information
%% NOTE: left commented -- filled out by the technical editor
%% \Volume{}
%% \Issue{}
%% \Month{}
%% \Year{}
%% \Submitdate{}
%% \Acceptdate{}
\Pages{1--xx}

\Address{
  Arquímedes De León Chacón Chacón\\
  iMetric\\
  Spain\\
  E-mail: \email{arquimedeschacon@gmail.com}
}

%% end of declarations %%%%%%%%%%%%%%%%%%%%%%%%%%%%%%%%%%%%%%%%%%%%%%%

\begin{document}

\section{Introducción} \label{sec:introduction}

% TODO: motivate the problem -- why measures of dependence beyond
% Pearson/Spearman/Kendall matter, why xi in particular, and the gap
% this paper fills (no comprehensive multi-factorial Monte Carlo study
% crossing n x distribution shape x variable type for xi, confirmed
% absent from the Austrian Journal of Statistics specifically).

\section{Marco teórico} \label{sec:framework}

% NOTA: seccion redactada en español para facilitar la revision.
% Traducir al ingles (y a title/sentence case segun estilo AJS)
% en la version final antes de enviar.

\subsection{El coeficiente xi de Chatterjee}

Las tres medidas clásicas de asociación --- $r$ de Pearson, $\rho$ de
Spearman y $\tau$ de Kendall --- son efectivas para detectar
dependencia lineal o monótona, pero pueden acercarse a cero incluso
ante una relación funcional perfecta, sin ruido, pero no monótona
entre $X$ e $Y$ \citep{chatterjee2021}. \cite{chatterjee2021} propuso
un coeficiente basado en rangos que supera esta restricción,
conservando la simplicidad y la tratabilidad asintótica de los
coeficientes clásicos.

Sean $(X_1,Y_1), \dots, (X_n,Y_n)$ copias i.i.d.\ de $(X,Y)$, con $Y$
no constante casi seguramente. Reordenamos los pares de modo que
$X_{(1)} \le \cdots \le X_{(n)}$ (los empates entre los $X_i$ se
rompen al azar), sea $r_i$ el rango de $Y_{(i)}$ entre todos los
$Y_j$, y sea $l_i$ el número de $j$ tales que $Y_j \ge Y_{(i)}$. El
coeficiente de correlación de Chatterjee se define como
\begin{equation} \label{eq:xi}
  \xi_n(X,Y) := 1 - \frac{n \sum_{i=1}^{n-1} |r_{i+1} - r_i|}
                         {2 \sum_{i=1}^{n} l_i (n - l_i)},
\end{equation}
que, en ausencia de empates entre los $Y_i$, se simplifica a
$\xi_n(X,Y) = 1 - 3\sum_{i=1}^{n-1}|r_{i+1}-r_i| / (n^2-1)$.

\cite{chatterjee2021} demuestra que $\xi_n(X,Y)$ converge casi
seguramente, cuando $n \to \infty$, a un límite determinístico
\begin{equation} \label{eq:xi-limit}
  \xi(X,Y) := \frac{\int \mathrm{Var}\big(E(\mathbf{1}_{\{Y \ge t\}} \mid X)\big)\, d\mu(t)}
                    {\int \mathrm{Var}\big(\mathbf{1}_{\{Y \ge t\}}\big)\, d\mu(t)},
\end{equation}
donde $\mu$ es la ley de $Y$. El límite satisface $\xi(X,Y) \in
[0,1]$, con $\xi(X,Y) = 0$ si y solo si $X$ e $Y$ son independientes,
y $\xi(X,Y) = 1$ si y solo si existe una función medible $f$ tal que
$Y = f(X)$ casi seguramente --- sin importar si $f$ es lineal,
monótona, u oscilante. No se requiere ningún supuesto distribucional
sobre $(X,Y)$, y $\xi_n$ es computable en tiempo $O(n \log n)$.

Tres propiedades de $\xi_n$ son directamente relevantes para el
diseño del estudio de simulación que sigue.

\emph{Direccionalidad.} A diferencia de los coeficientes de Pearson,
Spearman o Kendall, $\xi_n$ no es simétrico en $X$ e $Y$:
$\xi_n(X,Y) \ne \xi_n(Y,X)$ en general. Esto es intencional ---
$\xi_n(X,Y)$ cuantifica qué tan bien $Y$ está determinada por $X$, no
una asociación mutua y simétrica \citep{chatterjee2021}.

\emph{Sesgo de muestra finita.} Aunque $\xi_n$ converge a $\xi \in
[0,1]$, su valor máximo alcanzable para $n$ finito es estrictamente
menor que uno. \cite{dalitz2024} demuestran que, en ausencia de
empates entre los $Y_i$, $\xi_n(\vec x, \vec y) \le (n-2)/(n+1)$ para
toda permutación $\vec x$ del covariable --- para $n=20$ esta cota es
de solo $\approx 0.86$. Como esta cota es consecuencia de la
estructura algebraica de \eqref{eq:xi} y no del ruido en los datos,
garantiza un sesgo sistemático a la baja de $\xi_n$ como estimador de
$\xi$ cada vez que $\xi$ está cerca de uno, un efecto que debería
reducirse a medida que $n$ crece.

\emph{Varianza dependiente de la distribución bajo empates.} Bajo
independencia de $X$ e $Y$, $\sqrt{n}\,\xi_n(X,Y)$ es asintóticamente
normal con media cero \citep{chatterjee2021}. Si $Y$ es continua, la
varianza límite es una constante universal, $2/5$, libre de la ley de
$Y$ (Teorema 2.1). Si $Y$ \emph{no} es continua, en cambio, la
varianza límite $\tau^2$ depende de la ley de $Y$ (Teorema 2.2);
Chatterjee calcula $\tau^2 = 1$ para $Y \sim
\mathrm{Bernoulli}(1/2)$, exactamente 2.5 veces el valor del caso
continuo. Más aún, cuando los $X_i$ contienen empates, $\xi_n(X,Y)$
es en sí mismo un estadístico aleatorizado, ya que su valor depende
de cómo se rompan los empates. Estas no son curiosidades empíricas
sino consecuencias demostradas de \eqref{eq:xi}, y afectan
directamente a las variables discretas y ordinales del tipo que se
analiza rutinariamente en ciencias del comportamiento y sociales.

\subsection{Medidas competidoras de dependencia no lineal}

Otras dos familias de estadísticos han perseguido históricamente el
mismo objetivo --- detectar dependencia más allá de la linealidad o
la monotonía --- por vías matemáticas distintas, y son los puntos de
comparación naturales frente a los cuales debe evaluarse $\xi_n$.

\emph{Distance correlation.} \cite{szekely2007} definieron la
covarianza de distancias y la distance correlation (dCor) a partir de
las distancias par a par dentro de cada muestra, mostrando que
dCor$(X,Y) = 0$ si y solo si $X$ e $Y$ son independientes, para
vectores aleatorios $X$ e $Y$ de dimensión arbitraria. dCor es
simétrica y multivariante por construcción, pero, al operar sobre
distancias crudas en vez de rangos, no es invariante ante
reparametrizaciones monótonas no lineales de las marginales de la
forma en que sí lo son $\xi_n$, $\rho$ de Spearman o $\tau$ de
Kendall. Su costo computacional ingenuo es $O(n^2)$ (existe un
algoritmo rápido $O(n \log n)$ solo para el caso univariante), y la
prueba de independencia generalmente requiere un test de permutación
en vez de una distribución nula tratable.

\emph{Medidas basadas en cópulas.} \cite{hoeffding1948} propuso una de
las primeras pruebas no paramétricas generales de independencia, un
$U$-estadístico basado en rangos (la $D$ de Hoeffding) que mide la
discrepancia entre la distribución empírica conjunta de $(X,Y)$ y el
producto de sus marginales. $D$ es simétrica y consistente frente a
alternativas amplias, pero tiene un punto ciego conocido: ciertas
configuraciones dependientes y simétricas (por ejemplo, relaciones
tipo $Y = |X|$) pueden llevar a $D$ cerca de cero pese a existir
dependencia genuina. \cite{bergsma2014} introdujeron $\tau^\ast$, un
estadístico relacionado de covarianza de signos diseñado para ser
estrictamente positivo siempre que $X$ e $Y$ sean dependientes,
corrigiendo esta deficiencia. Ambos estadísticos son
$U$-estadísticos de orden cuatro o cinco, con un costo computacional
ingenuo de $O(n^2)$ o peor.

La Tabla~\ref{tab:comparison} resume las propiedades que distinguen a
$\xi_n$ de estas dos familias y que motivan su inclusión como puntos
de comparación en el diseño de simulación (Sección~\ref{sec:design}).

\begin{table}[hbt]
\caption{\label{tab:comparison}Propiedades de $\xi_n$ y sus medidas de dependencia de referencia.}
\begin{center}
\small
\begin{tabular}{p{4.2cm}p{1.6cm}p{2.1cm}p{2.6cm}}
\toprule
Propiedad & $\xi_n$ & dCor & $D$ Hoeffding / $\tau^\ast$ \\
\midrule
Simétrico en $(X,Y)$ & No & Sí & Sí \\[2pt]
Basado en rangos / invariante a monótonas & Sí & No & Sí \\[2pt]
Soporte multivariante nativo & No\footnotemark & Sí & No \\[2pt]
Costo computacional & $O(n \log n)$ & $O(n^2)$\footnotemark & $O(n^2)$ o peor \\[2pt]
Distribución nula tratable & Sí (Normal) & No (permutación) & No (permutación) \\
\bottomrule
\end{tabular}
\end{center}
\end{table}
\footnotetext[1]{\cite{azadkia2021} extienden $\xi_n$ a entornos condicionales y multivariantes; este estudio se limita al coeficiente bivariante original.}
\footnotetext[2]{Un algoritmo rápido $O(n\log n)$ para dCor existe solo en el caso univariante.}

\subsection{Efectos hipotetizados del tamaño de muestra, tipo de variable, forma de la distribución, estructura y régimen de dependencia}

El diseño factorial de la Sección~\ref{sec:design} no es un barrido
sin guía de condiciones arbitrarias: cada factor se eligió porque la
estructura algebraica o asintótica de $\xi_n$ revisada arriba predice,
o al menos permite, un efecto específico sobre su comportamiento.

\begin{description}
\item[H1 (tamaño de muestra).] Como la cota superior de muestra finita
  $(n-2)/(n+1)$ (Sección 2.1) se aproxima a uno solo cuando
  $n \to \infty$, $\xi_n$ debería exhibir un sesgo sistemático a la
  baja para $n$ pequeño que se reduce a una tasa algebraica conocida
  --- no simplemente la vaga expectativa de que ``más datos ayudan''.
  Además, este sesgo debería ser más pronunciado cuanto más alto sea
  el valor poblacional real $\xi(X,Y)$, ya que la cota solo restringe
  de forma apreciable los valores cercanos a uno.
\item[H2 (tipo de variable).] Como la varianza asintótica de $\xi_n$
  bajo independencia es libre de distribución ($2/5$) solo cuando $Y$
  es continua, y como los empates entre los $X_i$ vuelven a $\xi_n$ un
  estadístico aleatorizado, se espera que las variables ordinales y
  discretas (por ejemplo, ítems tipo Likert) alteren tanto la varianza
  como la reproducibilidad de $\xi_n$ respecto a las variables
  continuas --- una propiedad demostrada del estimador, no una
  preocupación especulativa.
\item[H3 (forma de la distribución).] La invariancia de $\xi_n$ ante
  transformaciones estrictamente monótonas de las marginales sugeriría,
  en principio, estabilidad frente a la forma de la distribución. Sin
  embargo, \cite{ansari2026} demuestran que $\xi$ \emph{no} es continua
  respecto a convergencia débil en general, y que su continuidad solo
  ha sido establecida para familias específicas (elípticas,
  $\ell_1$-simétricas). \cite{bucher2025} llevan esto al extremo:
  construyen explícitamente (vía ``shuffles of min'') secuencias que
  convergen débilmente a la independencia mientras $\xi$ permanece fijo
  en cualquier valor arbitrario, incluyendo uno, con la consecuencia de
  que los tests de independencia basados en $\xi_n$ pueden tener
  potencia trivial y de que no existen intervalos de confianza
  significativos de nivel exacto o uniforme para $\xi$. Estos
  resultados son de \emph{peor caso}, obtenidos con construcciones
  adversariales que quedan fuera del alcance de este estudio. Nuestra
  hipótesis, en cambio, es de \emph{caso típico}: dado que la
  continuidad de $\xi$ solo se ha verificado formalmente para un
  conjunto acotado de familias de distribución, no está probado que
  esa estabilidad se extienda a las formas no adversariales --- pero
  comunes en ciencias del comportamiento --- que se incluyen en nuestro
  diseño factorial (asimetría, colas pesadas, mixturas), en contraste
  con la deriva numérica esperada en dCor bajo las mismas
  transformaciones. Cabe precisar que \cite{ansari2026} ya resuelven
  \emph{analíticamente}, con fórmula cerrada, el caso de las
  distribuciones elípticas (incluyendo la normal y la $t$ de Student;
  ver su Sección 5) --- por lo que esas dos formas marginales no
  constituyen, en sí mismas, un vacío de conocimiento, sino que se
  incluyen en nuestro diseño como \emph{validación}: si nuestra
  simulación no reproduce su valor exacto de $\xi(Y,X)$ para esos
  casos, ello señalaría un error en el código de simulación. El vacío
  genuino de H3 recae específicamente en las formas marginales
  \emph{no} elípticas ni $\ell_1$-simétricas de nuestro diseño
  (exponencial asimétrica, mixtura bimodal), para las cuales ninguna
  garantía analítica de continuidad existe todavía.
\item[H4 (estructura de dependencia).] \cite{chatterjee2021} y
  \cite{gao2024} ya reportan que $\xi_n$ es más potente que dCor para
  relaciones oscilatorias y no monótonas, pero menos potente para
  relaciones lineales, circulares o heterocedásticas. No se ha
  examinado si este trade-off persiste de manera uniforme al variar
  conjuntamente el tamaño de muestra y el tipo de variable.
\item[H5 (régimen de dependencia).] \cite{auddy2024} demuestran que el
  comportamiento de $\xi_n$ bajo alternativas locales depende
  cualitativamente de si el valor poblacional real $\xi_0$ está en cero
  o es positivo: para testear independencia pura ($\xi_0 = 0$), el
  umbral de detección de $\xi_n$ es subóptimo, del orden $n^{-1/4}$,
  frente al óptimo paramétrico $n^{-1/2}$; para testear grado de
  asociación ($\xi_0 > 0$), en cambio, $\xi_n$ resulta minimax óptimo
  con la tasa correcta $n^{-1/2}$. Se espera, por lo tanto, que la
  potencia y el sesgo de $\xi_n$ no dependan solo de $n$, la forma de
  la distribución, o el tipo de variable de manera uniforme, sino que
  el régimen de dependencia real (cercano a independencia vs.
  moderado/alto) module todos estos efectos --- por lo que la magnitud
  de la dependencia simulada debe tratarse como un factor cruzado
  explícito, y no solo como los dos extremos ``independencia'' vs.
  ``dependencia fuerte''.
\end{description}

\section{Diseño de simulación} \label{sec:design}

% NOTA: seccion redactada en español para facilitar la revision.
% Traducir al ingles en la version final antes de enviar.

\subsection{Diseño factorial}

El diseño cruza cinco factores, cada uno directamente motivado por una
de las hipótesis de la Sección~\ref{sec:framework}:

\begin{description}
\item[Tamaño de muestra] --- $n \in \{20, 50, 100, 500, 2000\}$ (H1).
\item[Tipo de variable] --- continua, binaria, ordinal (Likert de 3, 5 y 7 puntos), obtenidas discretizando una misma variable latente continua antes de calcular cada estadístico (H2).
\item[Forma marginal] --- normal, uniforme, exponencial (asimétrica), $t$ de colas pesadas, mixtura bimodal (H3).
\item[Estructura de dependencia] --- forma funcional: lineal, cuadrática, senoidal, circular, heterocedástica (H4).
\item[Régimen de dependencia] --- magnitud poblacional de $\xi$, aproximadamente $0, 0.1, 0.3, 0.5, 0.8$, calibrado mediante el nivel de ruido añadido a cada relación funcional (H5).
\end{description}

Para cada combinación de condiciones se calculan $\xi_n$, dCor, y
$D$ de Hoeffding / $\tau^\ast$ como puntos de comparación
(Sección~\ref{sec:framework}, Tabla~\ref{tab:comparison}).

\subsection{Mecanismos generadores de datos}

\subsubsection{Estructura de dependencia (H4)}

Para una marginal de $X$ dada (H3) y un nivel de ruido $\sigma \ge 0$,
generamos $Y$ según el modelo aditivo $Y = f(X) + \sigma E$, con $E$
independiente de $X$ y de varianza unitaria (su distribución sigue la
misma familia asignada a la marginal de $X$ en esa condición, salvo
para la estructura circular y heterocedástica, ver abajo). Las cinco
formas funcionales $f$ son:

\begin{align*}
\text{Lineal:} &\quad f(X) = X, \\
\text{Cuadrática:} &\quad f(X) = X^2, \\
\text{Senoidal:} &\quad f(X) = \sin(2\pi X), \\
\text{Circular:} &\quad Y = Z\sqrt{1-X^2} + \sigma E, \quad Z = \pm 1 \text{ con prob. } 1/2 \text{ (independiente de $X$)}, \\
\text{Heterocedástica:} &\quad Y = 3\big(s(X)(1-\lambda)+\lambda\big)E, \quad s(X) = \mathbf{1}(|X|\le 0.5),
\end{align*}
con $X$ reescalada a $[-1,1]$ en todos los casos. Las formas circular y
heterocedástica siguen las alternativas de \citet[Sección 4.3]{chatterjee2021},
usadas allí precisamente para exhibir condiciones donde $\xi_n$ es
\emph{menos} potente que dCor (H4); las tres primeras siguen a
\citet{dalitz2024} y \citet{gao2024}.

\subsubsection{Calibración del régimen de dependencia (H5)}

El nivel de ruido $\sigma$ no tiene el mismo significado para $\xi$ en
cada estructura funcional --- un mismo $\sigma$ puede corresponder a
valores muy distintos de $\xi(X,Y)$ según la forma de $f$. Por ello,
para cada combinación de (estructura, forma marginal, tipo de
variable), $\sigma$ se calibra numéricamente, no se fija a priori: se
aproxima el valor poblacional $\xi(X,Y)$ mediante Monte Carlo de gran
tamaño de muestra ($n=10^5$, promediado sobre múltiples réplicas;
usando la fórmula cerrada de \citet{ansari2026} como atajo y
verificación donde esté disponible, i.e., para las formas elípticas de
H3), y se ajusta $\sigma$ (búsqueda numérica, p.\ ej.\ bisección) hasta
que ese valor aproximado quede dentro de una tolerancia fija del nivel
objetivo $\xi_0 \in \{0, 0.1, 0.3, 0.5, 0.8\}$. El valor de $\sigma$
resultante para cada celda del diseño se documenta en el repositorio
de reproducibilidad (Sección~\ref{sec:design}, 3.4) junto con el
script de calibración.

\subsubsection{Discretización para tipo de variable (H2)}

Partiendo de la variable continua generada según lo anterior, las
versiones binaria y ordinales (Likert-3, Likert-5, Likert-7) se
obtienen aplicando la función de distribución acumulada teórica de la
marginal correspondiente para cortar el rango en $c \in \{2,3,5,7\}$
categorías de igual probabilidad, asignando códigos enteros
$1,\dots,c$. Esto preserva la estructura de dependencia subyacente
mientras introduce empates de forma controlada --- el mecanismo
exacto que H2 busca poner a prueba. La discretización se aplica,
según la condición, a $Y$, a $X$, o a ambos, para poder distinguir el
efecto de empates en el covariable (que aleatoriza $\xi_n$ directamente,
Sección~\ref{sec:framework}) del efecto de empates en la respuesta
(que altera $\tau^2$).

\subsection{Métricas de resultado}

Para cada condición, replicada $R$ veces:

\begin{description}
\item[Sesgo y MSE] de cada estadístico respecto a su valor poblacional teórico (H1, H3).
\item[Varianza empírica] entre replicaciones, incluyendo variabilidad atribuible al desempate aleatorio de rangos cuando aplica (H2).
\item[Error de tipo I y potencia] de la prueba de independencia asociada a cada estadístico, evaluados por separado en el régimen cercano a independencia y en el régimen de dependencia moderada/alta (H5).
\item[Tiempo de cómputo], dado que la ventaja de velocidad de $\xi_n$ es una de sus propiedades más citadas (Sección~\ref{sec:framework}, Tabla~\ref{tab:comparison}).
\end{description}

\subsection{Software y reproducibilidad}

% TODO: version de R, paquetes (XICOR, energy, etc.), politica de
% semilla fija, y estructura del repositorio de reproducibilidad
% (siguiendo el mismo patron ya usado en las validaciones Monte Carlo
% de AssumptionsLab: script + resultados + README por cada corrida).

\section{Resultados} \label{sec:results}

% NOTA: seccion redactada en español para facilitar la revision.
% Traducir al ingles en la version final antes de enviar.
%
% Cada tabla resume el efecto de UN factor (el de su hipotesis),
% manteniendo los otros cuatro fijos en un nivel de referencia
% (indicado en cada caption). La matriz factorial completa (todas
% las combinaciones de los 5 factores) queda como material
% suplementario/apendice reproducible.

\subsection{H1: efecto del tamaño de muestra}

\begin{table}[hbt]
\caption{\label{tab:h1}Sesgo de muestra finita de $\xi_n$ en función de $n$. Referencia fija: forma = normal, tipo de variable = continua, estructura = lineal, régimen = moderado ($\xi_0 \approx 0.5$).}
\begin{center}
\begin{tabular}{lccccc}
\toprule
 & $n=20$ & $n=50$ & $n=100$ & $n=500$ & $n=2000$ \\
\midrule
Media $\hat\xi_n$ &  &  &  &  &  \\
Sesgo             &  &  &  &  &  \\
MSE               &  &  &  &  &  \\
\bottomrule
\end{tabular}
\end{center}
\end{table}

\subsection{H2: efecto del tipo de variable}

\begin{table}[hbt]
\caption{\label{tab:h2}Varianza empírica y potencia de $\xi_n$ según tipo de variable. Referencia fija: $n$ de referencia, forma = normal, estructura = lineal, régimen = moderado.}
\begin{center}
\begin{tabular}{lccccc}
\toprule
 & Continua & Binaria & Likert-3 & Likert-5 & Likert-7 \\
\midrule
Varianza empírica &  &  &  &  &  \\
Potencia          &  &  &  &  &  \\
\bottomrule
\end{tabular}
\end{center}
\end{table}

\subsection{H3: estabilidad ante la forma de la distribución}

\begin{table}[hbt]
\caption{\label{tab:h3}$\xi_n$ frente a dCor bajo distintas formas marginales (misma dependencia subyacente). Referencia fija: $n$ de referencia, tipo de variable = continua, estructura = lineal, régimen = moderado. Las columnas marcadas (V) son de \emph{validación} contra el valor exacto de \citet{ansari2026} para familias elípticas (Ec.~33 y Sección 5.1); el vacío genuino recae en las columnas restantes (formas no elípticas, sin garantía analítica de continuidad).}
\begin{center}
\begin{tabular}{lccccc}
\toprule
 & Normal (V) & Uniforme & Exponencial & $t$, colas pesadas (V) & Mixtura bimodal \\
\midrule
$\xi_n$ &  &  &  &  &  \\
dCor    &  &  &  &  &  \\
\bottomrule
\end{tabular}
\end{center}
\end{table}

\subsection{H4: potencia según estructura de dependencia}

\begin{table}[hbt]
\caption{\label{tab:h4}Potencia de detección por estadístico y estructura de dependencia. Referencia fija: $n$ de referencia, forma = normal, tipo de variable = continua, régimen = moderado.}
\begin{center}
\begin{tabular}{lccccc}
\toprule
 & Lineal & Cuadrática & Senoidal & Circular & Heterocedástica \\
\midrule
$\xi_n$          &  &  &  &  &  \\
dCor             &  &  &  &  &  \\
$D$ / $\tau^\ast$ &  &  &  &  &  \\
\bottomrule
\end{tabular}
\end{center}
\end{table}

\subsection{H5: interacción entre régimen de dependencia y tamaño de muestra}

A diferencia de H1-H4, esta hipótesis es sobre una \emph{interacción}
--- el propio efecto de $n$ depende del régimen --- por lo que su
tabla es bidimensional en vez de un barrido de un solo factor.

\begin{table}[hbt]
\caption{\label{tab:h5}Potencia de $\xi_n$ cruzando tamaño de muestra y régimen de dependencia poblacional $\xi_0$. Referencia fija: forma = normal, tipo de variable = continua, estructura = lineal.}
\begin{center}
\begin{tabular}{lccccc}
\toprule
 & $\xi_0\approx 0$ & $\xi_0\approx 0.1$ & $\xi_0\approx 0.3$ & $\xi_0\approx 0.5$ & $\xi_0\approx 0.8$ \\
\midrule
$n=20$   &  &  &  &  &  \\
$n=50$   &  &  &  &  &  \\
$n=100$  &  &  &  &  &  \\
$n=500$  &  &  &  &  &  \\
$n=2000$ &  &  &  &  &  \\
\bottomrule
\end{tabular}
\end{center}
\end{table}

\section{Ilustración con datos reales} \label{sec:realdata}

% TODO: apply xi (and the competing measures) to the bfi dataset
% (Revelle, psych package, SAPA project): 2800 subjects, 25 Likert
% (1-6) personality items (Big Five), plus age (continuous), gender
% (binary), and education (ordinal, 1-5). Chosen because it spans
% all three variable types in the factorial design within a single
% real, non-toy, psychology dataset.

\section{Discusión} \label{sec:discussion}

\section{Conclusión} \label{sec:conclusion}

%\bibliographystyle{plainat}
\bibliography{paper}

\end{document}
